\begin{table}
\centering
\begin{small}
\begin{tabular}{|l|cc|cc|cc|c|}
\hline
platform & \multicolumn{2}{c|}{x86-64 CPU} & \multicolumn{2}{c|}{Arm64 CPU (M4)} & \multicolumn{2}{c|}{CUDA (T4)} & \multicolumn{1}{c|}{MPS (M4)} \\
library & JAX & PyTorch & JAX & PyTorch & JAX & PyTorch & PyTorch \\ \hline
atan2 & Inf (0.501)$^{\rm fn}$ & 0.501 & 256$^{\rm f}$ & 0.501 & 0.501 & 0.501 & 256 (0.501)$^{\rm f}$ \\
hypot & 180 (1.74)$^{\rm f}$ & 0.500 & 180 (1.74)$^{\rm f}$ & 0.500 & 0.501 & 0.500 & 2.37e21$^{\rm f}$ \\
pow & 2.98e8$^{\rm f}$ & 0.501 & 128$^{\rm f}$ & 0.501 & 0.501 & 0.501 & Inf (128)$^{\rm f}$ \\
\hline
\end{tabular}
\end{small}
\caption{bfloat16 bivariate: largest error over all pairs (exhaustive). Values: largest error in ulps under the rules of Gladman et al. (3 significant digits rounded up; an input where only one of the library result and the exact result is NaN gives Inf). In parentheses, when different: the largest error excluding results affected by a flush-to-zero/denormals-are-zero mode of the backend and NaN mismatches. 0.500 = correctly rounded on all inputs; -- = function not provided for this format/backend; marks: f = some results flushed by the backend's FTZ/DAZ mode, n = spurious NaN for some inputs, m = NaN missing (number returned where the result is NaN). }
\label{tab:bfloat16bi}
\end{table}
